\documentclass[conference]{IEEEtran}
\usepackage{fancyhdr}
\usepackage{graphicx}
\usepackage{listings}
\usepackage{xcolor}
\usepackage{amsmath}
\usepackage{hyperref}
\usepackage{indentfirst}
\usepackage{titlesec}
\usepackage{subcaption}
\setlength{\marginparwidth}{2cm}
\usepackage{todonotes}
\hypersetup{
    breaklinks=true,
    urlcolor=blue,
    linkcolor=black,
    citecolor=black
}

\lstset{
    language=C,
    basicstyle=\ttfamily\small,
    stringstyle=\color{red},
    commentstyle=\color{blue}{},
    breaklines=true
}

\titlespacing*{\section}
{0pt}                      % left margin
{2ex plus 1ex minus .2ex}  % space before
{1ex plus .2ex}            % space after

\titlespacing*{\subsection}
{0pt}
{3ex plus 1ex minus .2ex}
{0.8ex plus .2ex}

\title{Final Project Proposal: A Retrieval-Augmented Assistant for ASIC Design Debugging}
\author{
    \IEEEauthorblockN{Parker Schless (pis7)}
    \IEEEauthorblockA{ENMGT 5400: Applications of Artificial Intelligence for Engineering Managers \\
    Cornell University: School of Engineering Management}
}

\begin{document}
\maketitle

\section{Project Overview}
When debugging an ASIC design, the context you need to resolve a problem lives in your own project files: the RTL source code, the constraint files, the synthesis and place-and-route reports, and the DRC/LVS logs. These files are dense, structured, and not easy to search by hand, especially when a single design run can produce hundreds of pages of reports.

I propose to build and evaluate a retrieval-augmented generation (RAG) system that answers ASIC design questions by retrieving relevant context from design artifacts generated during my own processor tapeout work. Because commercial EDA tool documentation is distributed under non-disclosure agreements, the corpus will consist solely of design artifacts: RTL source files, constraint files, synthesis and place-and-route reports, timing reports, and DRC/LVS logs. To keep the project entirely on open infrastructure, I will run the full design flow on my Zeppelin processor targeting the Nangate 45nm open PDK, and use the resulting reports as the basis of the corpus. The idea behind RAG is straightforward: instead of relying on an LLM's training data alone, you first retrieve relevant documents from a knowledge base and then feed those documents to the model so it can give a grounded answer. The final deliverable will be the working system, a benchmark of 30--50 question--answer pairs drawn from real debugging episodes, and a written report evaluating how well the system works and where it falls short.

\section{Project Motivation}
I am in my final semester of the Cornell ECE MEng program, currently working on the Zeppelin superscalar processor tapeout, and I start full-time at Apple in physical design after graduation. For this project I will be running the Zeppelin design through the full synthesis, place-and-route, and STA flow targeting the open Nangate 45nm PDK, and using the resulting reports and logs as the corpus. Over the last several months I have spent a lot of time on exactly the kind of problems this tool would address: simulation failures after back-annotation, timing violations on clock-gating cells, disagreements between different EDA tools on the same design, and various DRC and LVS errors. In most of these cases the answer existed somewhere, whether in a user guide, a forum post, or a prior report in the project, but finding it was the slow part.

This project sits at the intersection of two things I want to get better at: the ASIC design flow itself, and building useful things with LLMs. I already have a corpus of real design files and debug logs that I know well, which means I can tell whether the system is giving me good answers or not. That is the hard part of evaluating any LLM-based tool, and having ground truth from my own experience is what makes this project feasible at this scale.

\section{Previous Work}
RAG was introduced by Lewis et al. (2020) as a way to improve LLM answers on knowledge-intensive tasks by retrieving relevant documents before generating a response\cite{lewis2020rag}. The basic pipeline has three steps: split your documents into chunks, convert those chunks into vector embeddings, and at query time retrieve the most similar chunks to use as context for the LLM. Since then, better embedding models and retrieval methods have made RAG a standard approach for building Q\&A systems over custom document collections.

The closest work to what I am proposing is ChipNeMo\cite{liu2023chipnemo}, where NVIDIA built RAG-based assistants for internal chip design tasks including EDA script generation and bug summarization. They showed that domain-adapted models outperform general-purpose ones on these tasks. My project is much smaller in scope: I am not fine-tuning models or working at NVIDIA's scale. Instead, I am applying the same basic RAG idea to a single student tapeout project and focusing on honestly evaluating how well it works, where it fails, and why.

\section{Promising Directions}
The main questions I want to explore are practical ones. First, how should I split the documents into chunks? Design reports are highly structured (timing path tables, violation summaries, netlist hierarchies), and naively chopping these into fixed-size pieces will break that structure and probably hurt retrieval quality. I want to compare a simple fixed-size chunking approach against one that tries to respect document structure.

Second, since the corpus is entirely design artifacts rather than a mix of documentation and design files, I want to understand how well the system can answer questions that would normally require tool documentation to resolve. This is a meaningful constraint: a question about why a timing violation is occurring can be partially answered from the reports, but may be harder to fully explain without the tool docs as context. Characterizing this limitation is itself a useful contribution.

Third, I want to figure out whether bad answers come from retrieving the wrong documents or from the LLM misinterpreting the right ones, since these are different problems with different fixes. I will build a benchmark where each question has a known answer and known source documents, so I can measure retrieval accuracy and answer quality separately.

\section{Schedule}
The project is scoped for the five weeks after spring break, roughly 30--35 hours of total effort.

\begin{table}[h]
\centering
\begin{tabular}{|c|l|p{5.5cm}|}
\hline
\textbf{Week} & \textbf{Date} & \textbf{Deliverable} \\
\hline
1 & Apr 13 & Corpus assembly: run the Zeppelin design through the full synthesis, place-and-route, and STA flow on Nangate 45nm; collect resulting RTL, constraints, reports, and DRC/LVS logs. Define document schema and metadata (design stage, tool, report type). \\
\hline
2 & Apr 20 & Baseline RAG pipeline: chunking, embedding model, vector store, and a simple query interface that returns an LLM-generated answer with source chunks. Create evaluation benchmark: 30--50 questions from my own debugging history, each with a known answer and supporting source documents. \\
\hline
4 & Apr 27 & Experiments: try different chunking strategies and retrieval configurations, measure how each change affects answer quality. Final write-up, error analysis, and short demo to be posted on YouTube or demonstrated in real time. \\
\hline
\end{tabular}
\end{table}

\section{AI Attribution}
I have used Claude for the initial explorations of this project by helping me find relevant papers on the topic, as well as understanding the basics of how RAG systems work. Claude also helped me formulate the problem statement and refine the project scope based on what is realistic for a five-week timeline. It also helped to write and proofread this proposal given specific constraints on what should be included. 

\bibliographystyle{IEEEtran}
\bibliography{references}

\end{document}